% STATUS: rewritten 2026-11-05
% STATUS: extended 2026-11-07 (Ateşli, Esen and Sütlü 2021 and Xu and Liu
% 2021 read at the abstract, the introduction and the main definitions)
% Sources read in full: Pojidaev and Saraiva 2006 (basic facts on M_8,
% pp. 3595--3596; inner derivations, p. 3596; Theorem 2.3, p. 3598), Kaygorodov
% 2014 St. Petersburg (the (n+1)-ary derivation definition, p. 576; Theorem 5,
% p. 582; Corollary 6, p. 583); Ateşli, Esen and Sütlü 2021 (abstract,
% pp. 1--2); Xu and Liu 2021 (abstract, pp. 1--2).
\section{n-Lie algebras and ternary algebras}\label{sec:nonassociative:n-ary}

The Hom-Lie algebras of the previous section keep a binary bracket. % claim-ok: own summary of the previous section

The classes of this section replace the bracket itself by an n-ary
operation. % claim-ok: own linking sentence to the n-ary classes

\begin{example}[{\cite[Section~1, pp.~3595--3596]{PojidaevSarahva2006}}]
\label{ex:nonassociative:m8}
Let $K$ be a field of characteristic different from $2$ and $3$. Let $A$
be a composition algebra over $K$ with involution $a\mapsto\bar a$ and
unit $1$, whose symmetric bilinear form
$\langle x,y\rangle=\tfrac{1}{2}(x\bar y+y\bar x)$ is nonsingular.
With the ternary multiplication
\[
\langle x,y,z\rangle
\;=\;(x\bar y)z-\langle y,z\rangle x+\langle x,z\rangle y
-\langle x,y\rangle z,
\]
$A$ becomes a ternary Malcev algebra, denoted $M(A)$. If $\dim A=8$, then
$M(A)$ is not a $3$-Lie algebra and is denoted by $M_8$.
\end{example}

Pojidaev and Saraiva call a derivation of $M_8$ inner if it belongs to the
Lie algebra generated by the right multiplications of $M_8$
\cite[Section~2, p.~3596]{PojidaevSarahva2006}.

\begin{theorem}[{\cite[Theorem~2.3, p.~3598]{PojidaevSarahva2006}}]
\label{thm:nonassociative:pojidaev}
All derivations of $M_8$ are inner.
\end{theorem}

The derivation algebra of $M_8$ is $21$-dimensional and isomorphic to the
simple Lie algebra $B_3$
\cite[Section~2, p.~3598]{PojidaevSarahva2006}. The next notion extends the
derivation concept itself to $n$-ary algebras. % claim-ok: own linking sentence

\begin{definition}[{\cite[p.~576]{Kaygorodov2014n1aryMaltsev}}]
\label{def:nonassociative:n1ary}
Let $A$ be an $n$-ary algebra. An \emph{$(n+1)$-ary derivation} of $A$ is a
collection $(f_0,f_1,\dots,f_n)$ of linear endomorphisms of $A$ such that
\[
f_0[x_1,\dots,x_n]\;=\;\sum_{i=1}^{n}[x_1,\dots,f_i(x_i),\dots,x_n]
\]
for arbitrary $x_1,\dots,x_n\in A$.
\end{definition}

An $(n+1)$-ary derivation $(f_0,f_1,\dots,f_n)$ with
$f_1=\dots=f_n$ is an $(n+1)$-ary quasiderivation, and the collections
arising from the centroid or from an ordinary derivation
$(D,D,\dots,D)$ are called trivial
\cite[pp.~576--577]{Kaygorodov2014n1aryMaltsev}. For the simple non-Lie
$n$-ary Mal'tsev algebras, nothing nontrivial occurs.

\begin{theorem}[{\cite[Theorem~5, p.~582]{Kaygorodov2014n1aryMaltsev}}]
\label{thm:nonassociative:kaygorodov-m8}
The simple ternary Mal'tsev algebra $M_8$ is free of nontrivial $4$-ary
derivations.
\end{theorem}

Kaygorodov also proves the same for the simple $7$-dimensional Mal'tsev
algebra $M_7$: it is free of nontrivial $3$-ary derivations
\cite[Theorem~1, p.~579]{Kaygorodov2014n1aryMaltsev}. Writing $\Delta(A)$
for the Lie algebra of generalized derivations of $A$ modulo its annihilator,
the two results combine as follows. % claim-ok: own linking sentence

\begin{corollary}[{\cite[Corollary~6, p.~583]{Kaygorodov2014n1aryMaltsev}}]
\label{cor:nonassociative:delta-m7-m8}
$\Delta(M_7)\cong\Delta(M_8)$.
\end{corollary}

The generalized derivations of an $n$-Lie algebra have a cohomological home. A
generalized derivation of an $n$-Lie algebra $L$ is a linear map
$D\colon L^{n-1}\to L$ that satisfies the two identities of Ateşli, Esen and
Sütlü \cite[equations~(2.9)--(2.10), p.~12]{AtesliEsenSutlu2021}. A derivation of
$L$ is the case of a map $D\colon L\to L$ with the $n$-ary Leibniz rule
\cite[p.~12]{AtesliEsenSutlu2021}.

\begin{theorem}[{\cite[Proposition~2.13 and Corollary~3.2, pp.~23 and~27]{AtesliEsenSutlu2021}}]\label{thm:nonassociative:n-lie-generalized-cocycle}
Let $L$ be an $n$-Lie algebra and let $D\colon L^{n-1}\to L$ be a generalized
derivation. Then the associated map $D^{\sharp}$ is a cocycle of the cohomology
of $L$ with coefficients in $\mathfrak{gl}(L)$, and the cohomology of the
generalized derivation extension $L\oplus k$ coincides with that of $L$.
\end{theorem}

The cohomology that classifies generalized derivations is not the usual
cohomology that classifies abelian extensions, and the two agree only for
$n=3$ \cite[p.~2]{AtesliEsenSutlu2021}. The paper also lifts the
Hochschild--Serre spectral sequence to $n$-Lie algebras
\cite[pp.~2--3]{AtesliEsenSutlu2021}.

A derivation can also be carried along as part of the structure. A 3-LieDer
pair is a 3-Lie algebra together with a derivation of it, and Xu and Liu define
a cohomology for such pairs with coefficients in a representation
\cite[p.~1]{XuLiu2021ThreeLieDer}. They use it to describe the deformations and
the abelian extensions of the pair \cite[p.~1]{XuLiu2021ThreeLieDer}. The same
paper raises the notion one level, to a 3-Lie$_2$Der pair, and shows that the skeletal
objects are classified by a triple \cite[p.~1]{XuLiu2021ThreeLieDer}.

This closes the tour of nonassociative classes. The next chapter returns to
associative and commutative algebras. % claim-ok: own closing link to the next chapter

Ebadian defines a quartic map $D$ of a ternary Banach algebra with $D([x,y,z])=[D(x),y,z]+[x,D(y),z]+[x,y,D(z)]$ as a fourth-order ternary derivation and proves Ulam stability generalizations of these derivations for the stated quartic functional equation on Fr\'echet algebras~\cite{WlUlamStabilityGeneralizationsOf2013}.

Dehghanian, Park and Sayyari introduce ternary antiderivations on ternary Banach algebras and prove their stability for the stated $(\alpha,\beta)$-functional inequality by the fixed point method~\cite{WlStabilityOfTernaryAntiderivati2023}.

Park and Lee revisit ternary quadratic derivations on ternary Banach algebras and $C^*$-ternary rings, noting that the earlier definition of Shagholi and others was not well defined, and prove their Hyers--Ulam stability results in the corrected setting~\cite{WlApproximateTernaryQuadraticDer2015}.
